\documentclass[../main.tex]{subfiles}
\graphicspath{{\subfix{../image/}}}

\begin{document}

\chapter{辛流形}\label{chapter:ch22-symplectic-manifolds}

在这最后一章中，我们在流形上引入一种新的几何结构，称为\emph{辛结构}；它表面上类似于 Riemann 度量，但结果却具有深刻得多的不同性质。它只不过是选取一个闭的、非退化的 $2$-形式。这个定义的动机起初可能并不明显，但随着我们看到这些性质如何被使用，它会逐渐显现出来。现在只需说明：非退化之所以重要，是因为它给出了切--余切同构，就像 Riemann 度量所提供的那样；而闭之所以重要，是因为它导出了光滑函数与流之间的深刻联系（见本章后面关于 Hamiltonian 向量场与 Poisson 括号的讨论）。辛结构在数学与物理中有着惊人广泛的应用，包括偏微分方程、微分拓扑和经典力学等众多领域。

在本章中，我们只能对辛几何这门学科作一个简短的概述。我们从讨论有限维向量空间上非退化交错 $2$-张量的代数开始，然后把注意力转向流形上的辛结构。最重要的例子，是在每个光滑流形的余切从上典范定义的辛结构。我们给出重要的 Darboux 定理的一个证明，它表明每个辛形式都可以通过选取光滑坐标而局部地化为典范形式；因此，与 Riemann 度量的情形不同，辛结构的“平坦性”没有局部障碍。

然后我们探讨辛结构最重要的应用之一。辛流形上任何光滑实值函数都导出一个典范的常微分方程组，称为 Hamiltonian 系统。这些系统是经典力学研究的核心。

在处理完 Hamiltonian 系统之后，我们对辛结构的一个奇维类似物——\emph{接触结构}——作简短介绍。然后在本章末尾，我们说明如何利用辛几何与接触几何来构造一阶偏微分方程的解。

\subfile{Chapter-22/section-01.tex}
\subfile{Chapter-22/section-02.tex}
\subfile{Chapter-22/section-03.tex}
\subfile{Chapter-22/section-04.tex}
\subfile{Chapter-22/section-05.tex}
\subfile{Chapter-22/section-06.tex}
\subfile{Chapter-22/section-07.tex}
\end{document}